\chapter{Giới thiệu}
\label{ch:gioi-thieu}

\section{Lý do chọn đề tài và bối cảnh nghiên cứu}
Trong kỷ nguyên số hóa hiện nay, đời sống sinh viên tại các trường Đại học đang đối mặt với nhiều biến động và thách thức to lớn. Sự chuyển dịch kinh tế kéo theo chi phí sinh hoạt, học phí và các nhu cầu dịch vụ học tập ngày càng gia tăng. Đối với phần lớn sinh viên, việc tự chủ cuộc sống khi bước vào môi trường đại học là một bước ngoặt lớn, đi kèm với áp lực tự quản lý bản thân trên nhiều phương diện. Trong số đó, hai vấn đề nổi cộm nhất ảnh hưởng trực tiếp đến chất lượng học tập và sự phát triển bền vững của sinh viên là quản lý tài chính cá nhân và cơ hội tiếp cận các nguồn lực hỗ trợ như học bổng học thuật.

Về khía cạnh quản lý tài chính cá nhân, sinh viên thường thiếu kiến thức nền tảng và kỹ năng quản lý dòng tiền. Thói quen chi tiêu tự phát, thiếu kế hoạch, và không kiểm soát được các khoản phát sinh dẫn đến tình trạng mất cân đối tài chính thường xuyên. Điều này không chỉ gây căng thẳng tâm lý mà còn gián tiếp làm suy giảm kết quả học tập khi sinh viên phải dành nhiều thời gian để đi làm thêm trang trải cuộc sống. Phương pháp quản lý tài chính "6 cái lọ" (6 Jars) của T. Harv Eker là một mô hình nổi tiếng thế giới, giúp phân bổ thu nhập một cách khoa học vào các quỹ chi tiêu thiết yếu, đầu tư, giáo dục, tiết kiệm dài hạn, hưởng thụ và từ thiện \cite{eker2005-6jars}. Tuy nhiên, việc áp dụng mô hình này một cách thủ công thường gặp nhiều khó khăn do tính phức tạp trong việc ghi chép, phân loại giao dịch hàng ngày và thiếu tính kỷ luật tự giác.

Về khía cạnh tiếp cận nguồn lực hỗ trợ, học bổng là một trong những mục tiêu quan trọng giúp sinh viên giảm bớt gánh nặng tài chính và tạo động lực phấn đấu trong học tập. Hiện nay, các nguồn thông tin về học bổng rất đa dạng nhưng phân tán từ nhiều nhà tài trợ, doanh nghiệp, tổ chức phi chính phủ đến chính sách nội bộ của từng trường. Sinh viên thường bỏ lỡ các cơ hội học bổng phù hợp do không tiếp cận được thông tin kịp thời hoặc hồ sơ năng lực cá nhân chưa được so khớp tối ưu với các tiêu chí xét tuyển phức tạp. Đồng thời, các nhà tài trợ cũng gặp khó khăn trong việc tìm kiếm đúng đối tượng sinh viên đáp ứng đầy đủ yêu cầu học thuật và hoàn cảnh cụ thể.

Xuất phát từ thực trạng trên, đề tài \textbf{"Hệ thống hỗ trợ sinh viên - phân hệ quản lý tài chính"} được nghiên cứu và phát triển. Đề tài hướng tới việc xây dựng một hệ sinh thái tích hợp đa nền tảng, áp dụng các giải pháp công nghệ hiện đại và trí tuệ nhân tạo (AI) để giải quyết triệt để hai bài toán cốt lõi: tự động hóa quản lý tài chính cá nhân theo mô hình 6 lọ và tối ưu hóa quy trình so khớp học bổng thông minh cho sinh viên.

\section{Mục tiêu và phạm vi nghiên cứu}
\subsection{Mục tiêu nghiên cứu}
Đề tài tập trung vào việc nghiên cứu, thiết kế và phát triển các mục tiêu cụ thể sau:
\begin{itemize}
    \item Nghiên cứu cơ sở lý thuyết và hiện thực hóa mô hình quản lý tài chính cá nhân "6 chiếc lọ" trên ứng dụng di động, giúp sinh viên thiết lập và duy trì thói quen chi tiêu khoa học.
    \item Ứng dụng trí tuệ nhân tạo (AI Agent và Large Language Models - LLM) để tự động hóa quy trình phân loại giao dịch tài chính từ ngôn ngữ tự nhiên, cũng như tư vấn tài chính cá nhân thời gian thực thông qua Chatbot trợ lý ảo.
    \item Xây dựng phân hệ học bổng trực tuyến cho phép nhà tài trợ tạo chương trình học bổng, cấu hình điều kiện và tài liệu yêu cầu; đồng thời hỗ trợ sinh viên lưu nháp, nộp, theo dõi và hủy hồ sơ ứng tuyển theo quy trình xử lý rõ ràng.
    \item Xây dựng hệ thống đẩy thông báo di động (FCM Push Notification) thông minh, chạy ngầm định kỳ nhằm đưa ra các cảnh báo tài chính chủ động và thông tin học bổng khẩn cấp tới sinh viên.
    \item Xây dựng phân hệ quản trị dữ liệu học thuật cho phép quản trị viên quản lý danh mục trường đại học, khoa/chương trình đào tạo và môn học; đồng thời hỗ trợ sinh viên lập bảng điểm cá nhân, theo dõi GPA, tổng tín chỉ và mục tiêu học tập.
\end{itemize}

\subsection{Phạm vi nghiên cứu}
\begin{itemize}
    \item \textbf{Về đối tượng nghiên cứu:} Tập trung vào hành vi tài chính, chi tiêu sinh hoạt thường nhật và hồ sơ năng lực học thuật của sinh viên đại học.
    \item \textbf{Về mặt công nghệ:} Phát triển hệ thống đa dịch vụ (Monorepo) tích hợp: Backend (NestJS, TypeScript), AI Service (FastAPI, Python, LangChain, Google Vertex API), Mobile App (Expo React Native), và Web Portal Admin (React JS, Material UI).
    \item \textbf{Về mặt dữ liệu:} Khai thác dữ liệu học thuật mẫu mô phỏng theo cấu trúc đào tạo tín chỉ tại các trường đại học lớn tại Việt Nam (điển hình là mô hình Đại học Quốc gia TP.HCM).
\end{itemize}

\section{Mô tả bài toán đề tài giải quyết}
\subsection{Bài toán quản lý tài chính cá nhân sinh viên theo phương pháp 6 lọ}
Bài toán yêu cầu xây dựng một mô hình số hóa 6 quỹ tài chính độc lập cho từng người dùng, đảm bảo các nguyên tắc:
\begin{itemize}
    \item \textbf{Tính phân bổ tự động:} Khi người dùng phát sinh nguồn thu nhập mới (lương, học bổng, trợ cấp), hệ thống phải tự động phân chia dòng tiền vào 6 lọ theo tỷ lệ phần trăm được cấu hình trước.
    \item \textbf{Tính toàn vẹn dữ liệu tài chính:} Số dư của từng hũ tài chính phải là kết quả tính toán động từ dòng tiền thu chi thực tế trong lịch sử giao dịch để tránh sai lệch số liệu. Khi có hành vi xóa hoặc rebalance (điều chỉnh tỷ lệ hũ), hệ thống phải thực hiện các giao dịch bù trừ nội bộ để đảm bảo tổng tài sản của người dùng không bị thay đổi bất thường.
    \item \textbf{Hạn mức ngân sách (Budgeting):} Cho phép người dùng giới hạn chi tiêu trong từng lọ theo chu kỳ tháng để cảnh báo khi chi tiêu vượt ngưỡng thiết yếu.
\end{itemize}

\subsection{Bài toán quản lý và ứng tuyển học bổng trực tuyến}
Bài toán đặt ra là làm thế nào để số hóa toàn bộ quy trình công bố, tiếp nhận và xét duyệt học bổng giữa nhà tài trợ và sinh viên. Mỗi chương trình học bổng có thông tin riêng về giá trị tài trợ, số lượng suất, thời hạn nộp, điều kiện GPA, ngành học, trường mục tiêu và bộ tài liệu yêu cầu. Hệ thống cần đảm bảo sinh viên có thể tiếp cận thông tin rõ ràng, chuẩn bị hồ sơ trực tuyến và theo dõi trạng thái xử lý sau khi nộp.

Đối với nhà tài trợ, hệ thống cần cung cấp cổng quản trị để tạo và cập nhật học bổng, cấu hình điều kiện xét tuyển, định nghĩa các tài liệu bắt buộc và xem danh sách hồ sơ ứng tuyển theo từng chương trình. Đối với sinh viên, hệ thống cần hỗ trợ lưu nháp, nộp hồ sơ, upload tài liệu minh chứng, cập nhật hồ sơ khi được yêu cầu bổ sung và hủy hồ sơ khi còn trong trạng thái cho phép.

\subsection{Bài toán quản trị cơ sở dữ liệu học thuật của trường học}
Để dữ liệu học thuật được sử dụng thống nhất giữa Webadmin và ứng dụng di động, hệ thống cần số hóa quy trình quản lý danh mục trường đại học, khoa/chương trình đào tạo và môn học. Mỗi trường có thể có nhiều chương trình đào tạo và nhiều môn học. Mỗi môn học cần lưu mã môn, tên môn, mã học phần, số tín chỉ, thang điểm và trạng thái hoạt động để phục vụ quá trình nhập điểm của sinh viên.

Đối với quản trị viên, hệ thống cần cung cấp giao diện tạo, cập nhật, xóa, tìm kiếm và import hàng loạt dữ liệu trường, khoa/chương trình đào tạo và môn học. Đối với sinh viên, hệ thống cần cho phép chọn trường và chương trình đào tạo từ danh mục có sẵn, sau đó sử dụng danh mục môn học của trường để nhập điểm, import bảng điểm, theo dõi GPA và tiến độ tín chỉ.

\subsection{Các câu hỏi nghiên cứu}
Để định hướng triển khai và đánh giá hiệu quả các giải pháp kỹ thuật đề xuất, đề tài tập trung trả lời hai câu hỏi nghiên cứu cốt lõi sau:
\begin{enumerate}
    \item \textbf{CH1 -- Phân loại giao dịch:} Làm thế nào để tự động hóa việc phân loại các giao dịch tài chính từ ngôn ngữ tự nhiên phi chuẩn (viết tắt, sai chính tả tiếng Việt) của sinh viên vào đúng hũ tài chính trong mô hình 6 lọ với độ chính xác cao mà không làm chậm trải nghiệm người dùng?
    \item \textbf{CH2 -- Quản lý học bổng:} Làm thế nào để thiết kế quy trình học bổng trực tuyến đảm bảo nhà tài trợ có thể công bố chương trình, sinh viên nộp hồ sơ hợp lệ, và hệ thống theo dõi được toàn bộ trạng thái xét duyệt từ lưu nháp, đã nộp, đang xét duyệt đến phê duyệt hoặc từ chối?
    \item \textbf{CH3 -- Hệ thống thông báo hiệu năng cao:} Kiến trúc hàng đợi tin nhắn cần được thiết kế như thế nào để đảm bảo tính sẵn sàng cao và không gây nghẽn hệ thống khi gửi đồng loạt hàng nghìn thông báo đẩy di động (FCM) tới các sinh viên cùng lúc?
\end{enumerate}

\section{Những thách thức và hướng giải quyết của đề tài}
\subsection{Những thách thức lớn}
Trong quá trình xây dựng hệ thống, nhóm nghiên cứu đối mặt với các thách thức chính sau:
\begin{enumerate}
    \item \textbf{Độ trễ và chi phí tích hợp AI:} Việc gọi trực tiếp các API của mô hình ngôn ngữ lớn (như Gemini) có độ trễ phản hồi tương đối cao và tốn kém tài nguyên.
    \item \textbf{Tính chính xác trong phân loại giao dịch:} Ngôn ngữ tự nhiên mà người dùng nhập vào khi ghi chép chi tiêu thường mang tính viết tắt, không chuẩn hóa ngữ pháp (ví dụ: "an trua 50k", "cf 30k").
    \item \textbf{Đồng bộ trạng thái giao diện thời gian thực:} Giao diện điện thoại di động yêu cầu cập nhật tức thì số dư của các hũ và biểu đồ báo cáo ngay sau khi người dùng thực hiện giao dịch mà không làm chậm trải nghiệm ứng dụng.
    \item \textbf{Bảo mật dữ liệu và phân quyền hệ thống:} Dữ liệu tài chính cá nhân và hồ sơ học thuật/học bổng của sinh viên là các thông tin nhạy cảm. Hệ thống cần đảm bảo an toàn truy cập dữ liệu, chỉ cho phép các đối tượng thích hợp tiếp cận các tập tài nguyên tương ứng (Sinh viên truy cập app cá nhân, Nhà tài trợ quản lý hồ sơ ứng tuyển, Admin quản trị hệ thống).
    \item \textbf{Đảm bảo tính nhất quán và toàn vẹn của dữ liệu giao dịch tài chính:} Khi thực hiện các tác vụ phức tạp như phân bổ thu nhập (income distribution) vào 6 hũ hoặc cập nhật thông tin hũ, việc phát sinh lỗi giữa chừng hoặc race conditions có thể dẫn đến sai lệch số dư.
    \item \textbf{Đồng bộ dữ liệu học thuật và bảng điểm cá nhân:} Danh mục trường, khoa/chương trình đào tạo và môn học cần nhất quán giữa Webadmin và Mobile App. Đồng thời, điểm sinh viên nhập phải được kiểm tra theo thang điểm của từng môn, tránh sai lệch GPA và tổng tín chỉ.
\end{enumerate}

\subsection{Hướng tiếp cận giải quyết}
Để vượt qua các thách thức trên, đề tài đề xuất các hướng giải quyết kỹ thuật cụ thể:
\begin{itemize}
    \item \textbf{Sử dụng kiến trúc AI Agent và Caching:} FastAPI AI Service được thiết kế độc lập, đóng vai trò điều phối (orchestration) sử dụng các Tool-calling Agent của LangChain để tối ưu hóa việc gọi mô hình. Áp dụng cơ chế lưu đệm Redis Caching với TTL (Time-to-Live) từ 5-15 phút tại Backend để lưu trữ các nhận định tài chính (AI Insights), giảm tải số lượng request trùng lặp gửi đến Gemini API.
    \item \textbf{Huấn luyện LLM Classifier với Rule Engine:} Sử dụng các bộ quy tắc gán cứng (Rule-based) kết hợp với kỹ thuật Few-shot Prompting để hướng dẫn mô hình Gemini phân loại các thuật ngữ viết tắt của người dùng một cách chính xác vào 6 lọ tài chính.
    \item \textbf{Đồng bộ giao diện qua React Query:} Áp dụng thư viện quản lý state React Query (TanStack Query) trên Expo Frontend để quản lý cơ chế Invalidation. Khi có bất kỳ thay đổi nào từ phía dữ liệu, hệ thống tự động làm mới bộ đệm cục bộ và render lại màn hình mà không cần reload toàn bộ app.
    \item \textbf{Xác thực JWT và cơ chế phân quyền RolesGuard:} Sử dụng JSON Web Token (JWT) để xác thực người dùng trên các API. Triển khai bộ phân quyền `RolesGuard` chặt chẽ trong NestJS để phân tách quyền truy cập dựa trên vai trò (Student, Donor, Admin), bảo vệ an toàn cho các endpoint nhạy cảm.
    \item \textbf{Quản lý dữ liệu qua Database Transactions và cơ chế recalculate số dư:} Áp dụng cơ chế transaction của TypeORM trong NestJS để đảm bảo tính nguyên tử (Atomicity) cho các tác vụ cập nhật tài chính và hồ sơ học bổng. Đồng thời, xây dựng hàm tính toán lại số dư (`recalculateJarBalance`) dựa trên lịch sử giao dịch thực tế để chống sai lệch.
    \item \textbf{Thiết kế cơ sở dữ liệu quan hệ chuẩn hóa cao trên PostgreSQL:} Xây dựng cấu trúc cơ sở dữ liệu với các ràng buộc khóa ngoại (Foreign Keys) chặt chẽ tại Webadmin để đồng bộ dữ liệu học thuật phân cấp từ trường học tới sinh viên, đảm bảo tính nhất quán toàn vẹn dữ liệu.
\end{itemize}

\section{Các đóng góp của đề tài}
Đề tài mang lại các đóng góp thiết thực cả về mặt học thuật lẫn ứng dụng thực tiễn:
\begin{itemize}
    \item \textbf{Về mặt công nghệ:} Hiện thực hóa thành công một kiến trúc Monorepo đa dịch vụ ổn định, giải quyết bài toán giao tiếp bất đồng bộ và đồng bộ trạng thái giao diện thời gian thực giữa thiết bị di động, máy chủ API và lõi AI.
    \item \textbf{Về mặt ứng dụng tài chính:} Số hóa thành công mô hình quản lý tài chính 6 lọ, bổ sung thêm các tính năng thông minh như cảnh báo ngân sách tự động, gợi ý tái cấu trúc hũ hằng tháng (Jars Rebalancing) và trợ lý ảo tư vấn tài chính thân thiện với sinh viên.
    \item \textbf{Về hỗ trợ học tập:} Xây dựng phân hệ quản lý dữ liệu học thuật và bảng điểm cá nhân, giúp sinh viên theo dõi GPA, tín chỉ tích lũy, mục tiêu tín chỉ và quản lý điểm theo danh mục môn học chuẩn hóa từ nhà trường.
\end{itemize}
